% ============================================================
%  第十章  实际原因（10.1—10.2）
% ============================================================
\chapter{实际原因}

\begin{flushright}
\kaishu\small
如今只待\\
查明此结果的原因，\\
或者说，查明此缺陷的原因，\\
因为这有缺陷的结果必有其原因。\\[2pt]
——莎士比亚《哈姆雷特》II.ii.100--4\\[3pt]
{\itshape
And now remains\\
That we find out the cause of this effect,\\
Or rather say, the cause of this defect,\\
For this effect defective comes by cause.\par}
\end{flushright}

\section*{序言}
\addcontentsline{toc}{section}{序言}

本章对“实际原因”（actual cause）这一概念给出形式化的阐析。所谓实际原因，是指在特定情景中被判定为对某一给定结果之产生负有责任的事件，例如：“苏格拉底饮下毒堇汁，是苏格拉底之死的实际原因。”人类直觉在察觉与确认这类因果关系方面极为敏锐，因此它被视为构造解释的关键（\S 7.2.3），也被视为判定法律责任的终极标准（即所谓“事实上的原因”，cause in fact）。

然而，尽管实际因果性在自然思维中无处不在，它却不是一个容易形式化的概念。Wright (1988) 引入的一个典型例子考虑两场火向一座房子推进。若火 A 在火 B 之前烧毁了房子，我们（以及全国各地的许多陪审团）肯定会把火 A 视为损害的实际原因，尽管任何一场火单独都足以烧毁房子，而且没有一场是必要的。可见，实际因果性所需要的信息超出了必要性与充分性；在原因与结果之间起中介作用的实际过程必须纳入考虑。但在结构模型的语言中，“过程”究竟是什么？因果过程的哪些方面定义了实际因果性？我们又该如何把关于情景中不确定方面的证据拼合起来，从而计算实际因果性的概率？

本章提出一个可以在结构模型语义学中表述的实际因果性理论。该理论立足于将在 \S 10.2 定义的维持（sustenance）概念：它把必要性与充分性的各个方面结合起来，用以度量原因在模型发生某些结构变化时仍能维持结果的能力。我们通过例子说明，这一理论如何避开 Lewis (1986) 的反事实依赖理论所遇到的问题，以及它如何既可用于生成关于特定情景的解释，又可用于计算这类解释事实上正确的概率。

\section{引言：必要因果性的不足}

\subsection{单一事件原因再探}

相对于特定情景而作的“一场车祸是 Joe 死亡的原因”这类陈述，被归入“单称的”（singular）、“单一事件的”（single-event）或“个例层面”（token-level）的因果陈述。相对于某类事件或某类个体而作的“车祸会致人死亡”这类陈述，则被归入“一般的”（generic）或“类型层面”（type-level）的因果断言（见 \S 7.5.4）。我们把单一事件陈述中的原因称为\emph{实际原因}，把类型层面陈述中的原因称为\emph{一般原因}。

类型断言与个例断言之间的关系，在哲学文献中一直有争议（Woodward 1990; Hitchcock 1995）；而“何者在先”“一个层面能否化归为另一个层面”这类优先性问题（Cartwright 1989; Eells 1991; Hausman 1998），又把注意力从更基本的问题上引开了：“类型陈述与个例陈述对我们的世界作出了哪些可检验的断言？因果知识又是如何组织起来的，以支撑这类断言？”这一争论导致一些理论把类型断言与个例断言看作两类不同的因果关系（如 Good 1961, 1962），各自需要一套自己的哲学说明（例如见 Sober 1985; Eells 1991, chap.~6）——用 Hitchcock (1997) 的话说，“这实在不是一个令人愉快的处境”。与之相反，结构理论把类型断言与个例断言视为同一种关系下的实例，二者的差别仅在于就问题所引入的情景特定信息的细节。正因如此，结构理论为研究这两个层面之断言的解剖结构、支撑每个层面所需的信息、以及哲学家们何以觉得其间关系如此难以理清，提供了一个形式化的基础。

结构理论的基本构件是函数 $\{f_i\}$，它们表示似律的机制，并为类型层面与个例层面的断言同时提供信息。这些函数是类型层面的，因为它们表示变量之间一般的反事实关系，适用于每一个假设情景，而不限于已经实现的情景。同时，这些关系的任何具体实例化又表示一个个例层面的断言。把不同情景区分开来的要素，由背景变量 $U$ 表示。当所有这些因素都已知，即 $U=u$ 时，我们手中便有了一个“世界”（定义 7.1.8）——它是对某一特定情景的理想而完整的描述，其中一切相关细节都已写明，没有任何东西留给机遇或猜测。在世界层面作出的因果断言，将是个例因果断言的极端情形。然而一般来说，我们并不具备指明单一世界 $U=u$ 所需的详尽知识，于是用一个概率 $P(u)$ 来概括我们对这些细节的无知。这就把我们带到概率因果模型的层面（定义 7.1.6）。基于这类模型、不涉及实际情景而作出的因果断言，将被归入类型层面的断言。因果效应断言，例如
\[
P(Y_x=y)=p,
\]
就是这类断言的例子，因为它们表达的是 $x$ 引起 $y$ 的一般倾向——这一倾向是对一切潜在情景作判断而得出的。\footnote{偶尔，因果效应断言甚至可以基于不完整的概率模型作出，此时只给出 $G(M)$ 与 $P(y)$——这就是识别问题（第三章）。但仅凭这样的基础，若不对 $\{f_i\}$ 或 $P(u)$ 有所了解，就无法作出任何个例层面的陈述（这里当然假定 $x$ 与 $y$ 已知确实发生）。} 然而在大多数情形中，我们对手头的情景只掌握部分信息——例如，知道 Joe 死了，知道他遭遇了车祸，或许还知道他开的是跑车、头部受了伤。这类事件特定信息的总体称为“证据”（$e$），可用来把 $P(u)$ 更新为 $P(u\mid e)$。由模型 $\langle M,P(u\mid e)\rangle$ 导出的因果断言，表示深浅不一的个例断言，其深浅取决于 $e$ 的具体程度。

由此可见，在结构理论中，类型断言与个例断言之别是一个程度问题。我们收集的事件特定证据越多，就越接近个例断言与实际原因的理想。PS 与 PN 这两个概念（第九章的主题）代表这一谱系上的中间点。概率充分性（PS）接近类型层面的断言，因为它不考虑实际情景，事实上把实际情景排除在考虑之外。概率必要性（PN）则对实际情景有所涉及，尽管只是初步的涉及（即 $x$ 与 $y$ 为真）。本节将把更多信息纳入考虑，以求更接近实际原因这一概念。

\subsection{抢先与结构信息的作用}

在 \S 9.2 中我们提到，PN 与 PS 都是因果模型的全局（即输入—输出）特征，只依赖于函数 $Y_x(u)$，而与中介于原因（$x$）与结果（$y$）之间的过程结构无关。这样的结构在因果解释中确有作用，下面的例子即可见之。

考虑图 10.1 所示的一个由灯泡与两个开关组成的电路。从使用者的角度看，灯对两个开关的反应是对称的：任一开关都足以把灯打开。然而在其内部，开关 1 闭合时不仅接通灯，还把开关 2 从电路中断开，使其失效。于是，当两个开关都闭合时，我们会毫不犹豫地宣称开关 1 是灯泡中电流的“实际原因”，因为我们知道，在这一特定事态下，开关 2 对电路通路不可能有任何影响。PN 与 PS 之中没有任何东西能够解释这种不对称；二者都基于响应函数 $Y_x(u)$，因而对电路的内部运作一无所知。

\begin{figure}[htbp]
\centering
\begin{tikzpicture}[causal,scale=1.0,every node/.style={transform shape}]
  % 电池
  \draw (0,3.2)--(2.86,3.2);
  \draw (3.14,3.2)--(6.1,3.2);
  \draw[thick] (2.86,3.02)--(2.86,3.38);
  \draw[thick] (3.14,3.10)--(3.14,3.30);
  % 左侧：灯泡
  \draw (0,3.2)--(0,0.6)--(0.72,0.6);
  \draw (0.94,0.6) circle (0.20);
  \draw (0.80,0.46)--(1.08,0.74);
  \draw (0.80,0.74)--(1.08,0.46);
  \draw (1.16,0.6)--(1.6,0.6);
  \node[font=\small,below=0.8mm] at (0.8,0.36) {灯泡};
  % 上支路 → 开关 1 的上触点（闭合位）
  \draw (1.6,0.6)--(1.6,2.3)--(4.5,2.3);
  \node[circle,draw,fill=white,inner sep=0.6pt,minimum size=1.6mm] (c1) at (4.52,2.3) {};
  % 下支路 → 开关 2
  \draw (1.6,0.6)--(1.6,0.15)--(2.45,0.15);
  \node[circle,fill,inner sep=1.1pt] (p2) at (2.45,0.15) {};
  \draw[very thick] (2.45,0.15)--(3.0,0.62);
  \node[circle,draw,fill=white,inner sep=0.6pt,minimum size=1.6mm] (c2) at (3.02,0.64) {};
  \draw (3.02,0.64)--(4.2,0.64)--(4.5,1.33);
  \node[circle,draw,fill=white,inner sep=0.6pt,minimum size=1.6mm] (c3) at (4.52,1.35) {};
  \node[font=\small,below=1.4mm] at (2.7,0.02) {开关 2};
  % 开关 1
  \node[circle,fill,inner sep=1.1pt] (p1) at (5.2,1.62) {};
  \draw[very thick] (5.2,1.62)--(4.6,2.22);
  \draw (5.2,1.62)--(5.95,1.62)--(5.95,3.2);
  \node[font=\small,below=1.6mm] at (4.95,1.22) {开关 1};
  % ON/OFF
  \draw[{Stealth}-{Stealth}] (6.45,1.15)--(6.45,2.1);
  \node[font=\footnotesize,above=0.5mm] at (6.45,2.12) {开};
  \node[font=\footnotesize,below=0.5mm] at (6.45,1.13) {关};
\end{tikzpicture}
\caption{开关 1（而非开关 2）被知觉为灯亮的原因，尽管二者都不是必要的。}
\label{fig:10.1}
\end{figure}


这个例子代表了一类涉及\emph{抢先}（preemption）的反例，它们正是针对 Lewis 的反事实因果理论提出的。它说明了一个事件（例如开关 1 闭合）如何能被视为原因，尽管在其缺席时结果依然存在。Lewis (1986) 对此类反例的回答是修改反事实判据：只要 $x$ 与 $y$ 之间存在一条由中间变量构成的反事实依赖链，即链上每一环的输出都反事实地依赖于其输入，就可让 $x$ 成为 $y$ 的原因。对开关 2 而言，这样的链条并不存在，因为给定当前事态（即两个开关都闭合），把开关 2 打开或关上都不会使电路的任何部分（在电学上）受到影响。下面的例子可以更清楚地说明这一点。

\begin{example}[沙漠旅行者——仿 P.~Suppes]
沙漠旅行者 $T$ 有两个仇敌。仇敌 1 在 $T$ 的水壶里下毒；仇敌 2 不知仇敌 1 所为，开枪把水壶打穿，水漏光了。一周后，$T$ 被发现死亡，两个仇敌都供认了自己的行为与意图。陪审团必须判定：谁的行为才是 $T$ 之死的实际原因。
\end{example}

令 $x$ 与 $p$ 分别表示命题“仇敌 2 开枪”与“仇敌 1 在水里下毒”，并令 $y$ 表示“$T$ 已死”。除这些事件之外，我们还将使用中间变量 $C$（意指氰化物）与 $D$（意指脱水），如图 10.2 所示。函数 $f_i(pa_i,u)$ 并未在图 10.2 中显式画出，但按照对故事的通常理解，它们由图中各父变量的取值决定各子变量的取值：\footnote{为简便起见，本章以下略去“$\wedge$”符号。}
\begin{equation}
\begin{aligned}
c&=px',\\
d&=x,\\
y&=c\vee d.
\end{aligned}\tag{10.1}
\end{equation}
把 $c$ 与 $d$ 代入 $y$ 的表达式，我们得到一个简单的析取式
\begin{equation}
y=x\vee px'\equiv x\vee p,\tag{10.2}
\end{equation}
它在对称性上颇具欺骗性。

\begin{figure}[htbp]
\centering
\begin{tikzpicture}[causal,scale=1.12,every node/.style={transform shape}]
  \node[circle,fill,inner sep=1.2pt] (x) at (0,2.0) {}; \node[font=\small,above=1.2mm] at (0,2.0) {$x$};
  \node[circle,fill,inner sep=1.2pt] (p) at (4.6,2.0) {}; \node[font=\small,above=1.2mm] at (4.6,2.0) {$p$};
  \node[circle,fill,inner sep=1.2pt] (D) at (1.30,1.0) {}; \node[font=\small,above=1.0mm] at (1.30,1.0) {$D$};
  \node[circle,fill,inner sep=1.2pt] (C) at (3.30,1.0) {}; \node[font=\small,above=1.0mm] at (3.30,1.0) {$C$};
  \node[circle,fill,inner sep=1.2pt] (y) at (2.30,0.0) {}; \node[font=\small,above=1.0mm] at (2.30,0.0) {$y$};
  \draw[ed] (x)--(D); \draw[ed] (x)--(C); \draw[ed] (p)--(C);
  \draw[ed] (D)--(y); \draw[ed] (C)--(y);
  \node[font=\small,left=1.2mm,align=right] at (-0.14,2.0) {仇敌 2\\射击水壶};
  \node[font=\small,right=1.2mm,align=left] at (4.74,2.0) {仇敌 1\\投毒};
  \node[font=\small,left=1.4mm] at (1.16,1.0) {脱水};
  \node[font=\small,right=1.4mm] at (3.44,1.0) {氰化物摄入};
  \node[font=\small,right=1.2mm] at (2.44,0.0) {死亡};
\end{tikzpicture}
\caption{沙漠旅行者例子中的因果关系。}
\label{fig:10.2}
\end{figure}


这里我们以生动的符号看到了结构信息所起的作用。虽然 $x\vee x'p$ 确实与 $x\vee p$ 逻辑等价，但二者在结构上并不等价：$x\vee p$ 对于交换 $x$ 与 $p$ 是完全对称的，而 $x\vee x'p$ 告诉我们，当 $x$ 为真时，$p$ 全然没有影响——不仅对 $y$ 没有影响，而且对任何可能影响 $y$ 的中间条件都没有影响。正是这种不对称，使我们宣称 $x$ 而非 $p$ 是死亡的原因。

按 Lewis 的看法，$x$ 与 $p$ 的差别在于把各自与 $y$ 连接起来的链条之性质。从 $x$ 出发存在一条因果链 $x\to d\to y$，其中每一元素都反事实地依赖于其前件。从 $p$ 到 $y$ 则不存在这样的链条，因为当 $x$ 为真时，链条 $p\to c\to y$ 被抢先了（在 $c$ 处）；也就是说，无论 $p$ 如何，$c$ 都“卡”在假上。换个说法：虽然 $x$ 并不满足引起 $y$ 的反事实检验，但它的一个后果（$d$）满足；在 $x$ 与 $p$ 都为真的条件下，若不是因为 $d$，$y$ 就会为假。

Lewis 的链条判据保住了因果性与反事实之间的联系，但它相当特设；毕竟，对于一个像“实际原因”这样关键的概念——我们正是靠它来判定法庭上被告之有罪或无罪——为什么要把反事实依赖链的存在当作定义性的检验？基本的反事实判据确实体现了一种实用的理据：我们不愿为一项本来无法避免的损害惩罚某人，也希望鼓励人们留意那些其行动能造成实质差别的境况。然而，一旦行动与后果之间的反事实依赖被另一个原因的出现所破坏，再坚持连接二者的链条上的那些中间反事实依赖，又有什么益处呢？

\subsection{过度决定与准依赖}

Lewis 的链条还有一个问题：它未能刻画同时出现的析取式原因。例如，考虑图 9.1 中的行刑队，并设步枪手 A 与 B 同时开枪，杀死了囚犯。我们的直觉把每个步枪手都视为死亡的贡献性实际原因，尽管两个步枪手都不满足反事实检验，而且在对方在场时也都不支持一条反事实依赖链。

这个例子代表了一种称为\emph{过度决定}（overdetermination）的情形，它对反事实理论构成了严峻挑战。Lewis 对这一挑战的回答，是对反事实判据再作一次修补。他提出，反事实依赖链应被视为过程的内在属性（例如子弹从 A 飞向 D 的飞行），而由于特殊环境（例如子弹从 B 飞向 D 的飞行）导致的依赖消失，不应被视为依赖的内在丧失；我们仍应把这样的过程算作准依赖的——“只要环境有所不同”便如此（Lewis 1986, p.~206）。

Hall (2004) 指出，准依赖这一概念引出了一些棘手的问题：“第一，过程究竟是什么？第二，说一个过程在其内在品格上‘恰如’另一个过程，是什么意思？第三，我们究竟如何‘度量环境的多样性’？”我们将借助一个称为\emph{因果束}（causal beam）的对象（\S 10.3.1）来回答这些问题；可以把因果束看作对“过程”这一概念的结构—语义阐析。在简短考察 Mackie 的进路（它也处理实际因果性问题，只是视角不同）之后，我们将在 \S 10.2 回到链条与束，回到抢先与过度决定的问题。

\subsection{Mackie 的 INUS 条件}

上一节所遇到的问题，是哲学家们为使单一事件因果性（此处即“实际因果性”）获得满意的逻辑阐析所作的许多尝试中的典型。这些尝试似乎始于 Mill 的观察：没有哪个原因对其结果是真正充分或真正必要的（Mill 1843, p.~398）。此后提出的众多理论——基于对充分性与必要性条件更为精细的组合——都遭遇了无法克服的困难（Sosa and Tooley 1993, pp.~1--8）。Mackie (1965) 的论述似乎是在这一逻辑框架内对“实际因果性”给出半形式化阐析的最早尝试；他的方案被称为 INUS 条件，曾经极为流行。

INUS 条件是说：若 $C$ 是“一个本身对结果并非必要、但却充分的条件之中的一个不充分然而必要的部分”，则事件 $C$ 被觉察为事件 $E$ 的原因（Mackie 1965）。\footnote{这个缩略语中的两个否定与两个“但”，使 INUS 成为哲学文献中最不便于记忆的助记符之一。简化之后，它应读作：“一个充分条件集中的必要元素，NESS”（Wright 1988）。} 尽管为 INUS 给出精确表述的种种尝试（包括 Mackie 1980 本人的一些尝试）并未形成连贯的方案（Sosa and Tooley 1993, pp.~1--8），INUS 背后的基本想法却颇有吸引力：如果我们能把 $\{S_1,S_2,S_3,\dots\}$ 看作（对 $E$ 而言）一切极小充分条件集的汇集，那么当事件 $C$ 是某个 $S_i$ 的合取项时，$C$ 就是 $E$ 的 INUS 条件。进一步，若 $C$ 是 $E$ 的 INUS 条件，且在所处情形下 $C$ 对某个这样的 $S_i$ 而言是充分的，则 $C$ 被视为 $E$ 的原因。举例来说，若 $E$ 可写成析取范式
\[
E=AB\vee \bar C D,
\]
则 $C$ 之所以是 INUS 条件，是因为它是析取项 $\bar C D$ 的一个成员，而该析取项对 $E$ 是极小且充分的。于是，若在所讨论的场合 $D$ 出现，$C$ 就会被视为 $E$ 的原因。\footnote{Mackie (1965) 还要求 $E$ 的一切不以 $C$ 为合取项的析取项都不出现，但那样一来，Mackie 的定义就与 Lewis 的反事实检验完全相同了。此处我采用较宽泛的定义，以容纳同时原因与过度决定；见 Mackie (1980, pp.~43--7)。}

这一基本直觉为许多学科的研究者所共有。例如，法学者倡导一种称为 NESS 的关系（Wright 1988），它代表“充分集中的必要元素”，是以更简便的助记方式对 Mackie 的 INUS 条件的重述。在流行病学中，Rothman (1976) 提出了一个类似的准则——称为“充分成分”——用以认定某种暴露可被称为致病的原因：“我们说暴露 $E$ 引起疾病，若包含 $E$ 的一个充分原因是最先补全的那个充分原因”（Rothman and Greenland 1998, p.~53）。Hoover (1990, p.~218) 把 INUS 条件与计量经济学中的因果性联系起来：“在 Simon 的意义上引起另一变量的任何变量，都可被视为该另一变量的 INUS 条件。”

然而，所有这些方案都有一个根本缺陷：逻辑必然性与充分性的语言不足以阐析这些直觉（Kim 1971）。类似的结论隐含在 Cartwright (1989, pp.~25--34) 的分析之中：她起初为 INUS 的直觉所吸引，最终却不得不纠正 INUS 的错误。

逻辑方案的根本缺陷，源于它未能在句法上区分表示稳定机制（或用 Mackie 的术语，表示“倾向性关系”）的公式与表示情境条件的公式。这一局限最简单的表现见于换质位：“$A$ 蕴涵 $B$”逻辑等价于“非 $B$ 蕴涵非 $A$”，这就使非 $B$ 成为非 $A$ 的 INUS 原因。这与直觉相悖；由“疾病引起症状”，我们不能推出“消除症状将引起疾病消失”。换质位的失效还进一步带来转导（经由共同原因的推断）方面的麻烦：若疾病 $D$ 引起两个症状 $A$ 与 $B$，则治愈症状 $A$ 将（在 INUS 的逻辑方案中）蕴涵症状 $B$ 的消失。

另一组问题源于句法敏感性。设我们把 Mackie 的 INUS 条件用于图 9.1 的行刑队故事。若把囚犯死亡的条件写成
\[
D=A\vee B,
\]
则 $A$ 满足 INUS 判据，我们可以合理地断定 $A$ 是 $D$ 的原因。然而，若代入我们的模型中明确成立的 $A=C$，便得到
\[
D=C\vee B,
\]
于是 $A$ 突然不再作为合取项出现在 $D$ 的表达式中。难道我们要断定 $A$ 不是 $D$ 的原因吗？当然，我们可以禁止代入、坚持 $A$ 与 $B$、$C$ 一起留在析取式中，以此来避免这种消失。但这样一来，更糟的问题随之而来：在队长发出信号（$C$）而两个步枪手都未开枪的情形下，囚犯仍会被判定为死亡。简言之，传达故事中影响之流动的结构信息，无法用标准逻辑句法编码——Mackie、Rothman 与 Wright 的直觉必须重新表述。

最后，让我们考虑沙漠旅行者的例子。在那里，旅行者之死在式 \textup{(10.2)} 中表达为
\[
y=x\vee x'p.
\]
这一表达式并非极小析取范式，因为它可以改写为
\[
y=x\vee p,
\]
由此会得到一个反直觉的结果：$x$ 与 $p$ 是造成 $y$ 的平等伙伴。另一方面，若我们允许 $y=x\vee x'p$ 这样的非极小表达式，那么我们也大可以允许等价的表达式 $y=xp'\vee p$，由此又会荒谬地断定：只要有人打穿水壶（$x$），不往水里下毒（$p'$）就会成为我们这位旅行者不幸的原因。

现在我们回到结构分析，上述句法问题在其中不会出现。倾向性信息通过结构式或反事实式（例如 $v_i=f_i(pa_i,u)$，其中 $u$ 是一般的）传达，而情境信息通过指涉某一特定世界 $U=u$ 的命题式（例如 $X(u)=x$）传达。结构模型不允许任意的变换与代入，即使真值保持不变。例如，若把 $c$（氰化物摄入）理解为受一个独立于支配 $y$ 之机制的单独机制支配，那么在 $y=d\vee c$ 中代入 $c$ 的表达式就是不允许的。

借助结构分析，我们下面提出一个形式框架，以刻画 Mackie 与 Lewis 的直觉。我们的分析将立足于因果性的一个方面，称为\emph{维持}：它结合了充分性与必要性的要素，也把结构信息纳入考虑。

\section{生产、依赖与维持}

因果充分性的概率概念 PS（定义 9.2.2）提示了一条挽救反事实因果理论的途径。让我们再考虑行刑队例子中的对称过度决定。每个步枪手的射击都具有等于 1 的 PS 值（见 \textup{(9.43)}），因为在一个囚犯仍然活着的状态 $u'$ 中，每次射击都会导致囚犯死亡。这一高 PS 值支持我们的直觉：尽管 PN 值很低（$PN=0$），每次射击都是死亡的实际原因。于是，一个看似合理的论点是：我们的直觉对充分性有所考虑，因而可以用 PN 与 PS 分量的恰当配比，构造出实际因果性的适当判据。

Hall (2004) 表达了类似的期待。在分析反事实进路所面临的困难时，Hall 观察到存在两种因果性概念，反事实理论只刻画了其中之一；未能刻画第二种概念，很可能正是它与直觉冲突的原因。Hall 把第一种概念称为“依赖”，把第二种称为“生产”。在行刑队例子中，直觉认为每次射击都是死亡的平等“生产者”。相反，反事实理论只检验“依赖”，而它之所以失败，是因为囚犯的状态并不“依赖”于任何一次单独的射击。

依赖与生产这两个概念，分别与必要性和充分性密切平行。因此，我们对 PS 的表述很可能为 Hall 的生产概念提供形式基础，并成为把实际因果性形式化的一步。然而，这一纲领若要成功，必须先克服一个基本障碍：生产性因果性对情景特定信息视而不见（Pearl 1999），这从下面的考察可以看出。

因果性的依赖方面诉诸原因 $x$ 在面对某些列联（否则这些列联将否定 $y$）时维持结果 $y$ 的必要性（定义 9.2.1）：
\begin{equation}
X(u)=x,\qquad Y(u)=y,\qquad Y_{x'}(u)=y'.\tag{10.3}
\end{equation}
而生产方面则诉诸原因（$x$）在 $x$ 与 $y$ 皆不出现的情境（$u'$）中引起结果（$y$）的能力（定义 9.2.2）：
\begin{equation}
X(u')=x',\qquad Y(u')=y',\qquad Y_x(u')=y.\tag{10.4}
\end{equation}
比较这两个定义，我们注意到生产的一个奇特之处：要检验生产，我们必须暂时走出自己的世界，想象一个 $x$ 与 $y$ 皆不出现的新世界 $u'$，施加 $x$，再看 $y$ 是否出现。因此，“$x$ 生产了 $y$”这句话只能在其中 $x$ 与 $y$ 为假的世界 $u'$ 中为真；于是看起来：(a) 没有任何东西能够（通过考虑生产）解释实际世界中确实发生了的任何事件；(b) 关于实际世界 $u$ 所收集的证据，无法用于生产在其中得到定义的那个假设世界 $u'$。

为克服这一障碍，我们诉诸因果性的一个方面，称为\emph{维持}：它以生产的特征充实了依赖概念，同时又停留在 $x$ 与 $y$ 皆为真的世界 $u$ 之中。维持与依赖之别，在于 $x$ 被期望保护 $y$ 免受何种列联的侵扰。式 \textup{(10.3)} 中所考虑的列联是“情境性的”——即由一组专属于手头情景的环境 $U=u$ 演化而来——而我们现在坚持：$x$ 将维持 $y$，所抵御的是由模型本身的结构修改所演化出来的列联（Pearl 1998b）。

\begin{definition}[维持]
设 $W$ 为 $V$ 中的一组变量，$w,w'$ 是这些变量的具体取值。我们说，相对于 $W$ 中的列联，$x$ 在 $u$ 中\emph{因果地维持} $y$，当且仅当
\begin{equation}
\begin{aligned}
&\text{(i)}\quad X(u)=x;\\
&\text{(ii)}\quad Y(u)=y;\\
&\text{(iii)}\quad \text{对所有 } w,\ Y_{xw}(u)=y;\ \text{以及}\\
&\text{(iv)}\quad \text{对某个 } x'\neq x \text{ 与某个 } w',\ Y_{x'w'}(u)=y'\neq y.
\end{aligned}\tag{10.5}
\end{equation}
\end{definition}

式 \textup{(10.5)} 的维持特征体现在条件 (iii)，即 $Y_{xw}(u)=y$：它要求仅 $x$ 一项就足以维持 $y$。该条件读作：若在 $u$ 中把 $X$ 设为其实际值（$x$），那么即使把 $W$ 设为任何不同于实际值的值（$w$），$Y$ 在 $u$ 中仍将保持其实际值（$y$）。条件 (iv) 即 $Y_{x'w'}(u)=y'\neq y$，则把在这种不利条件下维持 $Y=y$ 的“责任”归于 $X=x$：若把 $X$ 设为别的值（$x'$），则 $Y$ 将在至少一种设定 $W=w'$ 下放弃其当前值（$y$）。合起来看，(iii) 与 (iv) 蕴涵：存在一种设定 $W=w'$，在其中 $x$ 对 $y$ 既必要又充分。

维持是对“实际原因”所施加的合理要求吗？让我们再看图 9.1 中致使囚犯死亡的那两颗子弹。我们认为 A 在此情景中是 D 的实际原因，因为即使没有 B，A 也会“单独”维持 D。但既然 B 确实发生了，我们又该如何在形式上表达我们的情景 $U=u$ 中 B 的缺席？如果我们希望在 $u$ 所创造的语境中（假设性地）压制 B，就必须使用一个结构列联，设想 B 被某种违反规则 $B=C$ 的外部干预（或“奇迹”）变为假；例如，步枪手 B 因某种机械故障而未能开枪。我们完全清楚这种故障并未发生，然而，正是通过把故事表示为图 9.1 那样的多阶段因果模型这一行为本身，我们就已承诺要考虑这类故障。

回想每个因果模型所代表的并非仅一个模型，而是一整组模型——对 do(.) 算子的每一种可能状态各有一个——考虑干预式列联便是每个此类模型的内在特征。换言之，模型中各机制的自主性意味着：每个机制都昭示着自己可能发生的故障，而这些故障标示出因果解释应当据以运作的列联。因此，把这类列联写入实际因果性（它是一种解释形式）的定义之中，是合理的。

定义 10.2.1 中 $W$ 的选取应当谨慎。我们显然不能允许 $W$ 包含一切中介于 $X$ 与 $Y$ 之间的变量，因为这样一来，任何 $x$ 都不可能维持 $y$。更严重的是，若不限制 $W$，我们就有取消真正抢先、把非原因变成原因的危险。例如，在沙漠旅行者故事（图 10.2）中取 $W=\{X\}$ 且 $w'=0$，我们就把仇敌 1 变成了死亡的实际原因，这既违背直觉，也违背实际情景（该情景排除了氰化物摄入）。提出“因果束”这一概念（Pearl 1998b），正是为了使 $W$ 的选取对实际情景的干扰最小。\footnote{Halpern and Pearl (1999) 允许选取任何集合 $W$，只要其补集 $Z=V\setminus W$ 被 $x$ 所维持，即对所有 $w$ 有 $Z_{xw}(u)=Z(u)$。}
